\chapter{Counting to Three}
\label{chap:finite-gauge-single-invariant}

\section{After the Floor}

The first chapter ended with a narrow permission. After the collapse of
truth, the grammar could not recover by announcing a universe of ready-made
objects. It could only sanction one local identity, a needle strong enough
to let a reading return to itself and weak enough not to smuggle in a whole
metaphysics. That needle is the hinge on which the present chapter turns.
With it, the grammar can stabilize truth long enough to ask what a finite
measurement must preserve when it is varied.

The answer is not yet a field, a force, or a law of nature. The answer is a
gauge: a finite rule for comparing neighboring readings while forbidding
the comparison from pretending to see beyond its own horizon. A gauge path
carries a sequence of permitted distinctions. An action assigns a cost to
that sequence. A variation tests whether an admissible slip changes the
reading. The whole construction remains finite. It may refine, but each
refinement must leave a record that the grammar can read.

This finite setting changes the meaning of stationarity. A stationary
reading is not a picture of something frozen. It is a path whose permitted
local slips leave no first residual to spend. The grammar reads the path,
reads the allowed variations, and finds that the residue vanishes. Because
the residue is the only authorized witness of imbalance, its vanishing is
the grammatical form of stationarity.

The chapter follows that single obligation through six renderings. First,
finite gauge paths and actions give stationarity its local reader. Second,
the vanishing first variation and the discrete Euler--Lagrange equations
become two ways to read one constraint. Third, a closed balanced cubic
spline shows that finite anchors can force the remaining residue to zero.
Fourth, real partitions and squeezed residuals explain how finite
truncations approach completion without pretending that the continuum was
given in advance. Fifth, a boundary anchor kills the nullspace that would
otherwise let unread activity pass as zero. Finally, the second variation is
identified as the unique positive invariant carried through the whole
chain.

The title of the chapter is therefore literal. A finite gauge is enough to
read one invariant, but only after the grammar has forbidden every cheaper
escape: unearned continuity, free boundary drift, invisible null modes, and
residual imbalance. What remains is not a decorative scalar. It is the
single quantity the grammar is forced to carry.

\section{Distinguishable}

A gauge path begins as a finite itinerary of distinguishable readings. It
does not require a completed space of possible histories. It requires only a
ledger of neighboring commitments: this reading can follow that one, this
slip is permitted, this endpoint is held fixed, this residue is still
unspent. The path is finite because the grammar has not earned any other
kind of record. It may name a direction of refinement, but it may not
collapse the direction into an already completed continuum.

The word gauge matters because the path is not merely a list. A list can
forget why its entries can be compared. A gauge path carries the rule by
which one entry is read against the next. It tells the grammar which local
changes are admissible, which labels are only names for the same reading,
and which changes would cross a boundary the current record cannot support.
The gauge therefore quotients away unearned presentation while preserving
the distinctions that still have measurement content.

An action is the finite cost of such a path. It is not an invisible
substance attached to the path, and it is not a completed integral waiting
above the finite record. It is the ledger's way of accumulating the cost of
admissible comparisons. If two adjacent readings agree in the required
sense, the cost reads as balanced there. If they fail to agree, the failure
appears as residue. The action is useful because it lets the grammar ask a
single question of many local comparisons: what remains after the permitted
balances have been paid?

The stationary reader is the rule that asks this question under variation.
A variation is not an arbitrary fantasy path. It is a permitted slip inside
the gauge. The endpoints that have been committed remain committed. The
labels that have been quotiented remain quotiented. Only the degrees of
freedom still open to the grammar may move. This restriction is not a
technical convenience. It is what prevents the reader from using a
variation to erase the very record it is supposed to test.

When a variation is applied, the action changes. The change can be read in
orders. The first order is the linear remainder, the part that would be
visible to an infinitesimal push in any admissible direction. But the
grammar does not need an infinitesimal as a new object. It needs the finite
discipline of comparison: shrink the permitted slip, read the leading
residue, and ask whether that residue can be forced to vanish for every
allowed test. If it can, the path is stationary.

This is why residual vanishing is the grammatical form of stationarity. A
path is not stationary because it looks smooth, short, elegant, or
physically plausible. It is stationary because the grammar has exhausted the
first readable imbalance. Every permitted local slip is tested, and no
linear residue remains. The stationary claim therefore has the shape:
within this gauge, with these endpoints, under these admissible variations,
the first residual vanishes.

A closed flat path gives the simplest picture. Suppose the path returns to
its boundary commitments without accumulating a readable slip. Each local
comparison is flat relative to the gauge, and the closing comparison
restores the initial commitment rather than adding a new debt. The grammar
then reads the whole loop as reducible: not because nothing happened, but
because every permitted difference has already been balanced by the
gauge. The residual ledger has no first-order entry left to spend.

Interpolation gives a useful illustration, provided it remains an
illustration. A thermometer scale may be fixed by two reproducible events
and then supplied with intermediate marks. The intermediate marks are not
new facts about the water. They are readings permitted by the rule that
relates the anchors. If a proposed curve between the anchors introduces a
distinguishable kink that refinement can detect, the kink appears as
residue. If the permitted completion removes every such first-order kink,
the completion reads as stationary relative to the chosen gauge.

The same discipline appears in a calibration path. A sequence of readings
can drift, but not every drift is an admissible variation. Some changes
only rename the scale; some change the relation between anchor events; some
expose a residual imbalance. The gauge identifies the harmless relabelings,
forbids the changes that abandon the committed anchors, and carries the
residue of the remaining changes. Stationarity is the moment when that
carried residue vanishes for all allowed slips.

This first section gives the chapter its central grammar. The finite gauge
permits paths, actions, and variations; it forbids unbounded escape from
the record; it identifies presentations that do not alter the reading; it
carries residuals until they vanish. Once this reader is available, the
next question is forced. If vanishing first residue is stationarity, what
local balance is being read when the residue vanishes?

\section{Admissible}

The stationary reader makes a demand before it supplies an equation. It
requires every admissible first residue to vanish. Classical notation often
writes this demand as
\[
  \delta A[\gamma;\eta] = 0
\]
for every permitted variation \(\eta\) of a path \(\gamma\). The notation is
helpful, but it is not the authority. The authority is the grammar that says
which variations are permitted, which endpoints remain fixed, and which
residue counts as first-order evidence of imbalance.

An endpoint is not held fixed because endpoints are sacred. It is held fixed
because the ledger has already committed it. If a variation could move the
committed boundary, the test would no longer be a test of the path between
the boundaries. It would be a new problem with a new record. The first
variation therefore reads only the interior degrees of freedom that the
gauge still permits to slip.

In the familiar variational calculation, the change of action separates
into an interior term and a boundary term. The boundary term vanishes when
the variation is zero at the endpoints. What remains is an integral or sum
against the arbitrary interior perturbation. If the first variation is zero
for every such perturbation, the coefficient of the perturbation must
vanish. The usual continuum expression is the Euler--Lagrange equation,
\[
  \frac{\partial L}{\partial q}
  - \frac{d}{dt}\frac{\partial L}{\partial \dot q} = 0 .
\]
The finite gauge reads the same obligation without requiring the continuum
to have been given first.

In the discrete reading, the path is a finite string of nodes. Each interior
node is touched from the left and from the right. A slip at that node changes
the cost of the left comparison and the cost of the right comparison. The
node is balanced exactly when those two changes cancel. The discrete
Euler--Lagrange equation is this nodewise balance. It says that no interior
node carries an unpaired first residual.

This is not a different theorem from stationarity. It is the same constraint
read at a different scale. The first-variation reading tests the whole path
against every admissible perturbation. The Euler--Lagrange reading inspects
the local balances that make all those tests vanish. If a node has a
nonzero residual, a variation supported near that node can read it. If every
supported variation reads zero, no node can retain a first residual. The
global vanishing and the local equations identify each other.

The grammar is doing important work here. It forbids a residue from hiding
behind a smooth phrase. A path may look continuous and still fail the local
balance. It also forbids the opposite mistake: a local equation may be
written down, but unless it is tied to the permitted variations and boundary
commitments, it has not yet been read as stationarity. The equation is not a
badge. It is the ledger entry left when the first variation has no remaining
debt.

The finite setting makes the equivalence especially sharp. Suppose an
interior node has a positive residual when pushed one way and a negative
residual when pushed the other. The path has not balanced that node. A
small supported variation exposes the imbalance immediately. Conversely, if
every interior node balances its incoming and outgoing costs, then every
finite supported variation decomposes into nodewise tests whose residues
sum to zero. The path is stationary because the grammar can find no first
residual anywhere it is allowed to look.

This also explains why the phrase ``least action'' must be handled with
care. A stationary path need not be globally least among every imaginable
path. The grammar has not licensed every imaginable path. It has licensed a
gauge, endpoints, and variations inside that gauge. The first claim is
therefore local and grammatical: no permitted first slip lowers or raises
the action to first order. Later positivity will say when this stationary
reading is also protected by a positive quadratic remainder.

The minimizing-variation example illustrates the distinction. A functional
with fixed endpoints is varied by a perturbation that vanishes at those
endpoints. Integration by parts moves the derivative off the perturbation
and leaves a term multiplying the free interior test. The classical
calculus says that this multiplier must vanish. The grammar says why that
move is legitimate: the boundary is pinned by commitment, the interior test
is arbitrary within the gauge, and a nonzero multiplier would be a readable
residue.

Thus the Euler--Lagrange equations are not imported as a law from outside
the finite grammar. They are the local reading forced by residual
vanishing. The finite gauge permits the test, carries the first variation,
identifies the local residuals, and forces their disappearance. Once that
has happened, the chapter can ask a more concrete question: what kind of
finite completion is sufficient to guarantee that no residual remains?

\section{Countable}

Finite measurement gives the grammar anchors before it gives the grammar a
curve. An anchor is a committed reading: an event, a mark, a boundary value,
a recorded position, a calibration point. Between anchors, the record is
silent. The silence is not empty permission. The grammar may complete the
record only in ways that remain answerable to the anchors and to the
residual tests already established. A completion that inserts a hidden
distinction has made a new claim, and the new claim must pay for itself.

A spline is the disciplined form of this completion. It does not say that a
continuous object was present behind the finite record all along. It says
that, once the anchors are fixed, the grammar may assign intermediate
values only by the least strained rule compatible with those anchors. The
intermediate values are consequences of the completion. They are not new
measurements. If a later refinement records a new anchor, the completion
must answer to it. Until then, the spline is the permitted smooth shadow of
the finite ledger.

The cubic case is the first case rich enough to carry the needed balance.
A line can connect two anchors but has no bending to account for. A
quadratic can bend but cannot generally balance neighboring pieces without
leaving a derivative debt at the knot. A cubic patch has enough freedom to
match value, slope, and bending moment across an anchor. It can carry the
position, the first slip, and the second slip without inventing a new
interior fact.

This is why the chapter speaks of a closed balanced cubic spline. Balanced
means that adjacent patches agree at their shared anchor in the readings the
gauge is allowed to test. The value agrees, so the path does not jump. The
first reading agrees, so the path does not hide a corner. The bending
reading agrees, so the path does not hide a curvature debt. Closed means
that the boundary of the chain has no unspent residue: the completion
returns to the committed boundary ledger without adding a final slip.

Once these requirements hold, the first residual has nowhere to live. Any
supported variation inside a patch is already governed by the cubic
balance. Any variation at a knot is answered by the matching conditions
from the two neighboring patches. Any variation at the boundary is
forbidden by the committed endpoint or cancelled by closedness. The grammar
therefore reads residue zero. Spline sufficiency is exactly this claim: a
closed balanced cubic completion is sufficient to force stationarity.

The claim is modest in the right way. It does not say that every stationary
path must be presented as a cubic spline. It says that this finite spline
grammar supplies enough structure to make residual vanishing unavoidable.
The anchors carry the finite facts; the cubic patches carry the least
completion; the knot balances carry the local Euler--Lagrange reading; the
closure carries the boundary debt to zero. Nothing else is needed for the
first residual to vanish.

The distinction between precision and accuracy is useful here. A spline can
make a value at an unmeasured point precise in the sense that the completion
identifies a unique value there. That precision is a grammatical
consequence of the anchors and the completion rule. It is not accuracy in
the empirical sense. If a later observation places a new anchor elsewhere,
the ledger changes and the completion must be recomputed. The grammar
therefore permits precise intermediate readings while forbidding them from
masquerading as independently recorded facts.

This is also why hidden curvature cannot be admitted for free. Suppose two
readers refine the same interval between the same anchors. If a hidden bend
were large enough to be distinguishable, refinement would expose a new
event. If no refinement exposes such an event, the permitted completion
must collapse toward the same balanced cubic patch. Agreement of value,
slope, and bending moment at shared anchors is not an aesthetic condition.
It is the grammar's way of forbidding a contradiction at the place where
two finite records meet.

Spline sufficiency therefore turns the stationary reader into a constructive
finite shape. The grammar has not assumed a continuum; it has permitted a
completion. It has not assumed smoothness; it has forced enough smoothness
to prevent unrecorded residue. It has not assumed a law of motion; it has
identified the closure condition under which all allowed first variations
read zero. The path is stationary because the balanced spline leaves no
first-order imbalance for the gauge to detect.

The next problem is scale. A single spline over a finite set of anchors
gives a local completion, but measurement also refines. Partitions split,
supports shrink, and residuals are tested on finer ledgers. The grammar
must explain how these finite completions are squeezed toward a completed
reading without treating the completed reading as primitive.

\section{Closure Under the Naturals}

The spline section shows how a finite set of anchors can force a balanced
completion. But a measurement grammar also has to live with refinement. A
partition can be split. A local support can be narrowed. A residual that was
invisible at one scale can become readable at another. If the chapter is to
reach a genuine invariant, it must explain how finite partitions approach a
completed reading without pretending that completion was available at the
start.

A real partition is a finite record laid across a real coordinate reading.
The word real does not license an infinite possession of every point. It
marks the kind of scale on which the finite record is being placed. The
grammar still carries only finitely many cells, finitely many anchors, and
finitely many supported tests. Each cell says: inside this bounded part of
the scale, these are the distinctions the ledger currently supports.

The finite truncation is therefore not a bad version of the continuum. It
is the only honest version the grammar can read at that stage. It carries
the supported data, the local completion rule, and the residual left by the
current mesh. A finer truncation may carry more data. It may split a cell,
add a knot, or test a narrower support. But it must preserve the older
commitments it has not explicitly revised. Refinement extends the ledger; it
does not give permission to forget why the older reading was admissible.

This is where the discrete Sobolev grammar enters. The grammar needs a way
to read not merely values, but activity: differences, bends, and supported
residuals. A Sobolev-style reading is the disciplined way of carrying both
the value information and the difference information. It measures whether a
finite record has bounded activity on its partition. It forbids a sequence
of refinements from hiding unbounded oscillation inside smaller and smaller
cells while pretending to converge as values.

The residual is then read weakly, against supported tests. A supported test
is a variation whose attention is confined to the part of the partition the
grammar has chosen. If every such test reads zero, the residual is invisible
on that supported ledger. If a test reads a nonzero value, the residual has
found a witness. The weak reading is weaker than omniscience and stronger
than silence: it asks exactly what the current partition is authorized to
ask.

Squeezing supplies the route to completion. At each finite stage, the
residual is bounded below by zero and above by a visible error bound. If the
upper bound is forced toward zero as the partition refines, the residual is
squeezed to zero. The completed reading is not asserted because someone has
looked at all points. It is read because every finite residual is trapped
between zero and a bound that vanishes under admissible refinement.

In symbols, the shape is simple:
\[
  0 \leq R_n \leq \epsilon_n,
  \qquad \epsilon_n \to 0 .
\]
The content is grammatical. The lower bound says no negative residue can be
used to cancel a hidden positive debt. The upper bound says the current
partition has measured the maximum unresolved imbalance it is still
carrying. The convergence of \(\epsilon_n\) says refinement is not wandering;
it is forcing the unread part to vanish. The limit is the reading permitted
by that forced vanishing.

The continuum-limit example is useful as a warning. Three coarse readings
can name a stress-strain relation, but they cannot determine every possible
continuum curve. A linear fit may converge to the wrong threshold if the
actual admissible shape contains a kink. Mesh refinement is honest only when
the residual bound being squeezed belongs to the grammar of the problem.
Refinement cannot rescue a model whose tests forbid the feature that the
record is trying to read.

This warning also protects the chapter from a common overreach. Completion
is not a license to add whatever smooth structure one likes. A completion is
admissible only when it is forced by finite partitions, supported tests, and
vanishing residual bounds. The Hilbert-style closure that later chapters
will use is therefore not a magical room added behind the finite record. It
is the completed home of the finite activity norm once the grammar has shown
that Cauchy refinements have a place to land.

Real partitions also explain why the spline result was not merely local
craftsmanship. A balanced cubic patch over one partition can be refined into
smaller patches. If the knot balances persist and the residual bound
shrinks, the finite completions form a tower. The tower carries values and
activity together. It reads convergence not by wishful smoothness but by the
disappearance of supported residuals.

Thus the grammar permits a continuum reading only at the end of a disciplined
squeeze. It starts with finite truncations, carries their supported
Sobolev activity, tests weak residuals, bounds the unread part, and forces
that bound to vanish. Completion is the quotient of all refinements that no
supported test can tell apart once the residual has vanished.

The next obstruction is subtler. A residual may vanish while a whole family
of invisible modes remains free. A constant drift, an affine tilt, or a
boundary-floating mode can carry no interior difference and yet still
prevent positivity. The grammar therefore needs an anchor.

\section{The Thirty-Six Step Ladder}

Residual vanishing is not yet positivity. A path may have no interior
imbalance and still drift as a whole. A chain may have zero differences
between neighboring nodes while every node has been shifted by the same
amount. A bending form may fail to see an affine tilt. These modes are not
errors in the calculation. They are null modes: differences the current
interior grammar cannot read.

The nullspace is the set of readings that cost nothing to the form being
tested. In a first-difference grammar, constant shifts vanish because every
neighboring difference is zero. In a bending grammar, affine motions may
vanish because the second difference is zero. If the only question is
interior balance, such modes are harmless. If the question is positivity,
they are fatal. A positive form must read zero only on the zero activity.

The anchor is the committed boundary reading that kills this nullspace. It
does not add arbitrary force to the system. It pins one part of the ledger
so that an otherwise invisible drift becomes readable. If every node moves
together but the pinned node is required to remain fixed, the drift is no
longer admissible. If an affine tilt would evade a bending reading but the
boundary fixes enough data to identify the tilt, the tilt collapses. The
anchor forces a null mode to confess as boundary motion.

This is why the anchor earns its keep. It is not a convenience placed after
the mathematics has finished. It is the condition under which the energy
reading becomes honest. Without the anchor, the form can say only that an
activity is zero modulo unread modes. With the anchor, the quotient is
resolved: the unread modes are forbidden by the boundary commitment, and
zero activity pins the state itself.

The grammar here recalls the needle from the previous chapter, but the
roles are different. The needle sanctioned one local identity after truth
collapsed. The anchor uses a committed boundary identity to prevent the
finite gauge from floating. The needle restored a place to stand. The
anchor makes that standing place expensive enough to distinguish drift from
rest. Both are bounded identity grants. Neither permits a universal sameness
to be smuggled in.

Boundary-value examples make the necessity plain. A differential relation
may be perfectly balanced and yet admit many solutions until boundary data
are supplied. A geometric reception problem may satisfy its propagation
relations and still have multiple admissible branches until an additional
event or environmental constraint collapses the ambiguity. The governing
relation carries an invariant, but uniqueness waits for a boundary
commitment. The anchor is that commitment in the finite gauge.

Once the nullspace is killed, positivity becomes readable. The quadratic
activity form can now say:
\[
  Q(u) \geq 0,
  \qquad Q(u)=0 \text{ only when } u=0 .
\]
The first inequality reads nonnegative activity. The second reads
definiteness. It is not enough for the form to avoid negative values; it
must also refuse to call a nonzero admissible activity zero. The anchor
supplies the refusal by forbidding the modes that the interior form cannot
see.

The phrase ``pinned boundary'' should therefore be read literally and
grammatically. Pinned means the boundary is not available for variation.
Boundary means the commitment is made where the interior record meets what
it must answer to. The pin is not an external hand holding a physical object
still. It is the rule that says which readings are already committed and
therefore cannot be spent as free variation.

This also clarifies why positivity belongs after stationarity, not before
it. Stationarity kills the first residual. Positivity protects the remaining
quadratic reading from degeneracy. If the first residual has not vanished,
the path is not balanced. If the quadratic reading is not positive after
anchoring, the path may be balanced but not measured by a genuine invariant.
The chapter needs both: vanishing first variation and positive second
variation.

The anchor closes the loophole left by real partitions. Refinement can
squeeze residuals toward zero, but if the tower carries an invisible drift,
the completion is not definite. The pinned boundary quotients out that
ambiguity by forbidding it. Once the nullspace vanishes, the activity norm
becomes a true measure of activity rather than a measure modulo drift. The
grammar can then read a single positive quantity.

That quantity is waiting in the second variation. The first variation has
vanished. The nullspace has been pinned. The squeezed residual has no hidden
place to hide. What remains is the quadratic remainder that measures
activity itself. The final section identifies that remainder as the chapter's
single invariant.

\section{The Single Invariant}

At a stationary path, the first variation has vanished. This does not mean
that nothing remains to be read. It means that the first readable imbalance
is gone. The next term in the action change is quadratic. It measures how
the action changes when the path is pushed in a permitted direction after
the linear debt has already been paid. This quadratic reading is the second
variation.

The second variation is not merely the next term in a formal expansion. It
is the first surviving measure of activity at stationarity. If the grammar
has already forced
\[
  \delta A[\gamma;\eta] = 0 ,
\]
then the leading nonzero response to a permitted variation has the form
\[
  \delta^2 A[\gamma;\eta,\eta].
\]
On the diagonal, where the same variation is read against itself, this is
the activity magnitude. It is the square of the energy reading in the sense
needed here: nonnegative, quadratic, and zero only when the anchored
activity itself has vanished.

The previous sections were necessary to make this sentence honest. The
stationary reader killed the first residual. The Euler--Lagrange balance
identified the local form of that killing. Spline sufficiency supplied a
finite completion that leaves no first-order knot debt. Real partitions
squeezed finite residuals toward completion. The anchor killed the nullspace
that would otherwise let nonzero activity read as zero. Only after those
obligations have been met can the second variation be called positive.

Positivity gives the invariant its force. A nonnegative quantity alone is
too weak; it may vanish on a whole quotient of unread modes. A stationary
condition alone is too weak; it may mark a saddle or a flat drift. A
completion alone is too weak; it may be smooth while carrying the wrong
activity. The anchored second variation combines the necessary features. It
is quadratic, it survives stationarity, it respects the finite completion,
and it is positive on every nonzero admissible activity.

Uniqueness follows by exclusion. A constant term cannot measure variation,
because it remains the same under every permitted slip. A linear term cannot
be the invariant at stationarity, because the grammar has forced it to
vanish. A null mode cannot be the invariant, because the anchor has pinned
it out of admissibility. A higher-order remainder cannot be the single
positive invariant, because it is not the first surviving measure and does
not supply the quadratic activity norm by which refinement is bounded. The
second variation is the only term left with the right grammar.

This is the title theorem of the chapter. The finite gauge does not merely
produce many possible quantities and then choose one by taste. It permits
finite paths, actions, variations, splines, partitions, residual bounds, and
anchors. It forbids unearned continuity, untested residue, and unpinned
null drift. Under those restrictions, the second variation is forced to be
the unique positive invariant.

The diagonal reading is especially important. A bilinear second variation
can compare two admissible variations, but the invariant appears when the
activity is compared with itself. The diagonal value is the energy square:
\[
  E(\eta)^2 = \delta^2 A[\gamma;\eta,\eta].
\]
The formula should be read grammatically. The left side names the magnitude
of admissible activity. The right side says how the finite gauge reads that
magnitude after first variation has vanished and the boundary has pinned
the nullspace. The equality identifies two renderings of one carried
quantity.

Once the invariant is available, many familiar physical readings can later
attach to it. A stiffness form can instantiate it. A Hilbert norm can be
opened from it. Charge, phase, and other labels can be tested against it.
But those later readings do not create the invariant. They inherit it. This
chapter's work is to show that the invariant is already forced by the
grammar of finite stationarity.

The single invariant is therefore not a metaphysical decoration placed on
top of measurement. It is what measurement must carry when it is finite,
gauge-bound, stationary, completed by admissible refinement, and anchored
against drift. It is the quantity that remains readable when every cheaper
reading has either vanished, been quotiented, been pinned, or been squeezed
to zero.

\section{What Counting Carries Forward}

The first chapter earned a place for truth to stand. This chapter has shown
what that standing place must carry when the grammar permits variation.
Finite gauge paths give measurement a way to move without abandoning its
record. Actions collect the cost of that movement. Stationarity reads the
vanishing of first residuals. Euler--Lagrange balance localizes that
vanishing. Cubic spline sufficiency gives finite anchors a balanced
completion. Real partitions squeeze supported residuals toward the completed
reading. The pinned boundary kills the nullspace. The second variation then
stands alone as the unique positive invariant.

The next chapter can therefore become concrete without changing the plot.
It will not need to invent a new invariant. It must supply a concrete
operator that carries the invariant already forced here. The finite gauge
has done its work: it has identified what must remain unchanged when
measurement is allowed to vary.
